% Ad Familiares 12.15 --- headnote
% ad-familiares-12-15, 29 May / 2 June 43 BC, Perga

\textit{P. Cornelius Lentulus Spinther the younger to the Senate, consuls,
praetors, tribunes, and Roman People, from Perga in Pamphylia ---
Perseus dateline \textit{Scr. Pergae 1-6 \S\S iiii K., \S{} 7 iiii Non. Iun.
a. 711 (43)}: the first six sections were composed on 29 May 43 BC; the
seventh, a postscript reporting Dolabella's repulse at Antioch, was added
on 2 June. The conventional date for the dispatch as filed is 2 June 43.
Lentulus was the son of P. Lentulus Spinther (consul 57, Cicero's recall),
serving in the East as proquaestor pro praetore; with Trebonius murdered
at Smyrna in January, the province of Asia was effectively in the hands
of Dolabella's adventurers until M. Brutus's neighbouring force in
Macedonia and Cassius's gathering army in Syria changed the balance.}

\textit{This is a formal senatorial despatch in the older style ---
salutation in full (\textit{S. V. L. V. V. B. E. V.}, ``if you and your
children are well, it is well; I and the army are well'') and tone
controlled throughout. Lentulus has to explain why he turned aside to
Rhodes and was humiliated there: the Rhodians, bound by a renewed treaty
to count the Senate's enemies as their own, instead shut his men out of
their harbour and water, kept open communications with Dolabella, and
appear to have stalled his approach long enough for Dolabella's lieutenants
to slip out of Lycia. He pleads the cause before their assembly, fails,
recovers the cargo-fleet at Lycia, pursues Dolabella as far as Sida, and
turns back to gather the province's revenues. Three days later news
arrives from deserters that Dolabella has been refused entry at Antioch
and is fleeing toward Laodicea with Cassius four days' march behind him;
Lentulus appends the postscript with measured confidence ``that this most
criminal bandit will pay his penalty sooner than expectation reckoned.''
The letter is invaluable as the contemporary report of the campaign that
ended in Dolabella's death at Laodicea later in the summer, and as a
study in the embarrassed but dutiful prose of a young senatorial officer
explaining a setback up the chain.}
